\documentclass[a4paper]{article}

\usepackage{graphicx}
\usepackage{amsmath,amssymb}
\usepackage{minted}
\usepackage{multirow}
\usepackage{float}
\usepackage{todonotes}
\usepackage{forest}
\usepackage{xstring}
\usepackage{xcolor}
\usepackage[most]{tcolorbox}
\usepackage{xparse}
\usepackage{natbib}

\usetikzlibrary{graphs,graphdrawing}
\usegdlibrary{layered}

% for external graphviz
\input{/home/donkarlo/Dropbox/repo/nd_language_ai_project/src/nd_language_ai/formal/computer/programming/tex/latex/code_clips/external_graphviz.tex}

% for tags
\input{/home/donkarlo/Dropbox/repo/nd_language_ai_project/src/nd_language_ai/formal/computer/programming/tex/latex/parts/control_sequence/kind/semtag/semtag.tex}

% for links
\input{/home/donkarlo/Dropbox/repo/nd_language_ai_project/src/nd_language_ai/formal/computer/programming/tex/latex/code_clips/seehere.tex}

\title{Kausal}
\date{}

\begin{document}
    \maketitle
    
    
    \section{Einleitung}
        
        Diese Wörter und Ausdrücke verwendet man, wenn man einen Grund, eine Ursache oder eine Begründung ausdrücken will. Sie beantworten oft die Frage: „Warum?“
    
    
    \section{Unterordnende Konjunktionen}
        
        \subsection{weil}
            
            \textbf{Grammatische Rolle:} unterordnende Konjunktion
            
            \textbf{Gebrauch:} \emph{weil} leitet einen Nebensatz ein. Das konjugierte Verb steht am Ende.
            
            \textbf{Beispiele:}
            \begin{itemize}
                \item Ich bleibe zu Hause, weil ich krank bin.
                \item Er lernt Deutsch, weil er in Österreich arbeitet.
            \end{itemize}
        
        \subsection{da}
            
            \textbf{Grammatische Rolle:} unterordnende Konjunktion
            
            \textbf{Gebrauch:} \emph{da} leitet ebenfalls einen Nebensatz ein. Das konjugierte Verb steht am Ende. Oft klingt \emph{da} etwas schriftlicher oder formeller als \emph{weil}.
            
            \textbf{Beispiele:}
            \begin{itemize}
                \item Da ich krank bin, bleibe ich zu Hause.
                \item Ich kann heute nicht kommen, da ich einen Termin habe.
            \end{itemize}
    
    
    \section{Beiordnende Konjunktion}
        
        \subsection{denn}
            
            \textbf{Grammatische Rolle:} beiordnende Konjunktion
            
            \textbf{Gebrauch:} \emph{denn} verbindet zwei Hauptsätze. Die normale Wortstellung bleibt erhalten.
            
            \textbf{Beispiele:}
            \begin{itemize}
                \item Ich bleibe zu Hause, denn ich bin krank.
                \item Er geht früh schlafen, denn er ist müde.
            \end{itemize}
    
    
    \section{Konjunktionaladverbien}
        
        \subsection{deshalb}
            
            \textbf{Grammatische Rolle:} Konjunktionaladverb
            
            \textbf{Gebrauch:} \emph{deshalb} steht im Hauptsatz. Meistens steht das Verb direkt danach.
            
            \textbf{Beispiele:}
            \begin{itemize}
                \item Ich bin krank, deshalb bleibe ich zu Hause.
                \item Er hat wenig Zeit, deshalb arbeitet er sehr schnell.
            \end{itemize}
        
        \subsection{deswegen}
            
            \textbf{Grammatische Rolle:} Konjunktionaladverb
            
            \textbf{Gebrauch:} \emph{deswegen} ist sehr ähnlich wie \emph{deshalb} und kommt oft in der gesprochenen Sprache vor.
            
            \textbf{Beispiele:}
            \begin{itemize}
                \item Ich habe Kopfschmerzen, deswegen gehe ich ins Bett.
                \item Sie war müde, deswegen ist sie früh nach Hause gegangen.
            \end{itemize}
        
        \subsection{daher}
            
            \textbf{Grammatische Rolle:} Konjunktionaladverb
            
            \textbf{Gebrauch:} \emph{daher} drückt eine Folge aus und wirkt oft etwas schriftlicher.
            
            \textbf{Beispiele:}
            \begin{itemize}
                \item Es hat die ganze Nacht geregnet, daher ist die Straße nass.
                \item Er hat lange geübt, daher spricht er jetzt besser Deutsch.
            \end{itemize}
        
        \subsection{darum}
            
            \textbf{Grammatische Rolle:} Konjunktionaladverb
            
            \textbf{Gebrauch:} \emph{darum} bedeutet hier „aus diesem Grund“.
            
            \textbf{Beispiele:}
            \begin{itemize}
                \item Ich habe morgen eine Prüfung, darum lerne ich heute Abend.
                \item Sie fühlt sich nicht gut, darum bleibt sie heute zu Hause.
            \end{itemize}
    
    
    \section{Präpositionen und feste Ausdrücke}
        
        \subsection{aus diesem Grund}
            
            \textbf{Grammatische Rolle:} fester Ausdruck
            
            \textbf{Gebrauch:} \emph{aus diesem Grund} ist eher formell und kommt oft in schriftlichen Texten vor.
            
            \textbf{Beispiele:}
            \begin{itemize}
                \item Er hat die Frist verpasst. Aus diesem Grund wurde der Antrag abgelehnt.
                \item Das Produkt war beschädigt. Aus diesem Grund habe ich es zurückgeschickt.
            \end{itemize}
        
        \subsection{wegen + Genitiv}
            
            \textbf{Grammatische Rolle:} Präposition
            
            \textbf{Gebrauch:} \emph{wegen} steht mit einem Nomen und drückt einen Grund aus.
            
            \textbf{Beispiele:}
            \begin{itemize}
                \item Wegen des schlechten Wetters bleiben wir zu Hause.
                \item Das Spiel wurde wegen eines technischen Problems abgesagt.
            \end{itemize}
        
        \subsection{aufgrund + Genitiv}
            
            \textbf{Grammatische Rolle:} Präposition
            
            \textbf{Gebrauch:} \emph{aufgrund} ist eher formell und häufig in offiziellen oder schriftlichen Texten.
            
            \textbf{Beispiele:}
            \begin{itemize}
                \item Aufgrund eines Unfalls kam es zu einer Verspätung.
                \item Aufgrund seiner Erfahrung bekam er die Stelle.
            \end{itemize}
        
        \subsection{aus ... Gründen}
            
            \textbf{Grammatische Rolle:} Präpositionalgruppe
            
            \textbf{Gebrauch:} Diese Struktur wird mit verschiedenen Adjektiven verwendet, zum Beispiel: \emph{aus gesundheitlichen Gründen}, \emph{aus privaten Gründen}.
            
            \textbf{Beispiele:}
            \begin{itemize}
                \item Aus gesundheitlichen Gründen konnte er nicht kommen.
                \item Sie hat aus privaten Gründen abgesagt.
            \end{itemize}
    
    
    \section{Kurzer Überblick}
        
        \begin{itemize}
            \item \textbf{weil, da}: unterordnende Konjunktionen; sie leiten einen Nebensatz ein, das Verb steht am Ende.
            \item \textbf{denn}: beiordnende Konjunktion; sie verbindet Hauptsätze, die Wortstellung bleibt normal.
            \item \textbf{deshalb, deswegen, daher, darum}: Konjunktionaladverbien; sie stehen im Hauptsatz.
            \item \textbf{wegen, aufgrund}: Präpositionen; sie stehen mit einem Nomen.
        \end{itemize}
    
    
    \section{Vergleich}
        
        \begin{itemize}
            \item Ich bleibe zu Hause, weil ich krank bin.
            \item Da ich krank bin, bleibe ich zu Hause.
            \item Ich bleibe zu Hause, denn ich bin krank.
            \item Ich bin krank, deshalb bleibe ich zu Hause.
            \item Wegen meiner Krankheit bleibe ich zu Hause.
        \end{itemize}
    
    %\bibliographystyle{plainnat}
    %\bibliography{/home/donkarlo/Dropbox/repo/nd_shared_research/bibliography/bibliography.bib}
\end{document}